% !TEX root = ../main.tex
\begin{ouverture}[grandsnombres]{11}{Statistiques et simulations}{Partie 4 · Probabilités et statistiques}
{Indicateurs\enspace·\enspace Échantillons\enspace·\enspace Moyennes observées\enspace·\enspace Monte-Carlo}
\objectif{Lire des graphiques et calculer des indicateurs statistiques.}
\objectif{Simuler une variable et la moyenne d’un échantillon.}
\objectif{Comparer des moyennes observées à une espérance.}
\objectif{Expérimenter un critère de dispersion et une estimation d’aire.}
\prerequis{Moyennes pondérées, fréquences et variables aléatoires.}
\prerequis{Listes et boucles Python ; le chapitre 12 sert de référence.}
\end{ouverture}
\essentiel
\essSection{1}{Décrire une série statistique}
\begin{definition}
Pour des valeurs $x_i$ d'effectifs $n_i$, l'effectif total est $N=\sum_i n_i$
et la fréquence de $x_i$ vaut $f_i=n_i/N$.
La moyenne est $\overline x=\sum_i f_ix_i$ et la variance descriptive est
$v=\sum_i f_i(x_i-\overline x)^2=\sum_i f_ix_i^2-\overline x^2$.
L'écart type descriptif vaut $\sqrt v$ ; on divise ici par $N$, pas par $N-1$.
\end{definition}
\begin{propriete}[Médiane et quartiles]
Dans une série ordonnée, la médiane est la valeur centrale si $N$ est impair,
et la moyenne des deux valeurs centrales si $N$ est pair.
On prend $Q_1$ au rang $\lceil N/4\rceil$ et $Q_3$ au rang $\lceil3N/4\rceil$.
L'étendue est le maximum moins le minimum ; l'écart interquartile vaut $Q_3-Q_1$.
\end{propriete}
\essSection{2}{Lire les bons graphiques}
\begin{savoirfaire}
Un diagramme en barres convient aux valeurs discrètes. Dans un histogramme,
l'aire de chaque rectangle représente l'effectif ou la fréquence ; pour des
classes de largeurs inégales, les hauteurs sont des densités.
Un diagramme en boîte représente minimum, $Q_1$, médiane, $Q_3$, maximum.
\end{savoirfaire}
\begin{center}
\begin{tikzpicture}[x=1cm,y=1cm]
\draw[otMuted,->](0,0)--(7.5,0);
\foreach \x in {1,2,3,4,5,6,7}{\draw[otMuted](\x,.08)--(\x,-.08)node[below,font=\small]{\x};}
\draw[otNavy,thick](1,.8)--(2,.8) (4,.8)--(7,.8);
\draw[otNavy,fill=otNavyTint](2,.4)rectangle(4,1.2);
\draw[otGreenDark,thick](3.5,.4)--(3.5,1.2);
\draw[otNavy](1,.5)--(1,1.1) (7,.5)--(7,1.1);
\end{tikzpicture}
\par\legende{Boîte de minimum 1, quartiles 2 et 4, médiane 3,5, maximum 7.}
\end{center}

\clearpage
\essSection{3}{Simuler une variable et une moyenne}
\begin{propriete}
En Python, \code{random()} renvoie un réel pseudo-aléatoire de $[0;1[$.
L'expression \code{random() < p} simule un succès de probabilité $p$.
Une fonction de simulation peut être répétée pour créer une liste de $n$ valeurs,
puis une moyenne observée \code{sum(L)/len(L)}.
\end{propriete}
\begin{savoirfaire}[Distinguer théorie et observation]
Calculer d'abord l'espérance $\mu$ et l'écart type $\sigma$ du modèle.
Simuler plusieurs échantillons de même taille et comparer leurs moyennes à $\mu$.
Une moyenne observée fluctue ; elle n'est pas obligatoirement égale à $\mu$.
Fixer une graine permet de reproduire une expérience informatique.
\end{savoirfaire}
\essSection{4}{Expérimenter avec de nombreux échantillons}
\begin{definition}[Un critère d'observation]
Pour $N$ échantillons indépendants de taille $n$, on peut compter la proportion
des moyennes $m$ vérifiant $|m-\mu|\leq2\sigma/\sqrt n$.
Le programme demande d'expérimenter ce critère : il ne donne pas une proportion
universelle exacte, et cette proportion n'est pas une probabilité de certitude.
\end{definition}
\begin{exemple}[Monte-Carlo]
Pour estimer une aire dans le carré unité, tirer des points $(x;y)$ uniformes
dans ce carré et compter la proportion de ceux qui appartiennent à la zone.
Pour le quart de disque $x^2+y^2\leq1$, multiplier cette proportion par 4
donne une estimation de $\pi$.
\end{exemple}
\begin{attention}
Un résultat de simulation dépend du tirage et de la taille de l'échantillon.
Une observation numérique, même répétée, ne démontre pas un théorème général.
\end{attention}

\exercices

\rubrique{Découverte}
\begin{blocLivre}{ex-11-1}
\begin{exercice}[1]{D’un tableau aux indicateurs}
Les valeurs 0, 1, 2 et 3 ont pour effectifs 2, 3, 4 et 1.\begin{questions}
\item Calculer l’effectif total et les fréquences.
\item Calculer la moyenne, la médiane et l’étendue.
\end{questions}
\end{exercice}
\end{blocLivre}
\begin{blocLivre}{ex-11-2}
\begin{exercice}[1]{Lire un diagramme en barres}
\begin{center}
\begin{tikzpicture}\begin{axis}[otLivreGraphe,width=96mm,height=40mm,xmin=8,xmax=18,ymin=0,ymax=5,xtick={10,12,14,16},ytick={1,2,3,4},xlabel={valeur},ylabel={effectif},axis lines=left,every axis x label/.style={at={(axis description cs:1,0)},anchor=north east,yshift=-4mm,font=\small,text=otMuted},every axis y label/.style={at={(axis description cs:0,1)},anchor=south west,yshift=1mm,font=\small,text=otMuted}]
\addplot[ycomb,very thick,otNavy,mark=*]coordinates{(10,2)(12,4)(14,3)(16,1)};
\end{axis}\end{tikzpicture}
\end{center}
\begin{questions}
\item Reconstituer les effectifs et calculer la moyenne.
\item Donner la médiane et les quartiles selon la convention du cours.
\end{questions}
\end{exercice}
\end{blocLivre}
\begin{blocLivre}{ex-11-3}
\begin{exercice}[1]{Tracer une boîte}
On donne la série ordonnée $1,2,2,3,4,4,5,7$.\begin{questions}
\item Déterminer la médiane et les deux quartiles.
\item Tracer une boîte et calculer l’écart interquartile.
\end{questions}
\end{exercice}
\end{blocLivre}
\begin{blocLivre}{ex-11-4}
\begin{exercice}[1]{Une fréquence observée}
Une variable de Bernoulli de paramètre $0{,}3$ a fourni l’échantillon $[0,1,0,0,1,1,0,0,0,1]$.\begin{questions}
\item Calculer sa moyenne observée.
\item Calculer l’écart absolu à l’espérance. Cet écart invalide-t-il le modèle ?
\end{questions}
\end{exercice}
\end{blocLivre}
\begin{blocLivre}{ex-11-5}
\begin{exercice}[1]{Comprendre un tirage simulé}
Une fonction Python renvoie 1 si \code{random() < 0.25}, sinon 0.\begin{questions}
\item Quelle loi cherche-t-elle à simuler ?
\item Quelles valeurs peut prendre la moyenne de quatre appels ?
\item La moyenne de quatre appels vaut-elle nécessairement $0{,}25$ ?
\end{questions}
\end{exercice}
\end{blocLivre}
\clearpage
\rubrique{Entraînement}
\begin{blocLivre}{ex-11-6}
\begin{exercice}[2]{Un histogramme de largeurs inégales}
Les classes $[0;10[$, $[10;30[$ et $[30;60]$ ont pour fréquences $0{,}2$, $0{,}5$ et $0{,}3$.\begin{questions}
\item Calculer les hauteurs d’un histogramme dont les aires représentent les fréquences.
\item Comparer les hauteurs et les fréquences. Pourquoi des hauteurs directement égales aux fréquences seraient-elles trompeuses ?
\end{questions}
\end{exercice}
\end{blocLivre}
\begin{blocLivre}{ex-11-7}
\begin{exercice}[2]{Transformer une série}
Une série $X$ a pour moyenne 12 et pour écart type 3. On définit $Y=2X+5$ et $Z=X-12$.\begin{questions}
\item Calculer les moyennes de $Y$ et de $Z$.
\item Calculer leurs variances et écarts types.
\item Interpréter le centrage $Z=X-12$.
\end{questions}
\end{exercice}
\end{blocLivre}
\begin{blocLivre}{ex-11-8}
\begin{exercice}[2]{Une aire sous une parabole}
Dans 400 tirages de points uniformes dans le carré $[0;1]^2$, 137 vérifient $y\leq x^2$. L’aire théorique, fournie ici, vaut $1/3$.\begin{questions}
\item Donner l’estimation de Monte-Carlo.
\item Calculer l’erreur absolue exacte.
\item Écrire le test Python qui reconnaît un point dans la zone.
\end{questions}
\label{figlivre-r11-08-zone-ex}
\figureLivre{r11-08-zone}

\end{exercice}
\end{blocLivre}
\begin{blocLivre}{ex-11-9}
\begin{exercice}[2]{Estimer le nombre pi}
Parmi 1000 points tirés uniformément dans le carré unité, 786 sont dans le quart de disque $x^2+y^2\leq1$.\begin{questions}
\item Donner l’estimation de $\pi$ et son erreur absolue à $10^{-5}$ près.
\item Expliquer le facteur 4.
\item Un second essai de 1000 points doit-il donner exactement le même résultat ?
\end{questions}
\label{figlivre-r11-09-zone-ex}
\figureLivre{r11-09-zone}

\end{exercice}
\end{blocLivre}
\begin{blocLivre}{ex-11-10}
\begin{exercice}[2]{Comparer des moyennes d'échantillons}
$X$ vaut 0 ou 2 avec probabilités égales. Quatre échantillons de taille 5 donnent les listes $[0,0,2,0,2]$, $[2,2,2,0,2]$, $[0,0,0,0,0]$, $[2,0,2,2,0]$.\begin{questions}
\item Calculer $\mu$ et $\sigma$ pour le modèle.
\item Calculer les quatre moyennes.
\item Quelle proportion satisfait $|m-\mu|\leq2\sigma/\sqrt5$ ?
\end{questions}
\end{exercice}
\end{blocLivre}
\begin{blocLivre}{ex-11-11}
\begin{exercice}[2]{Effet de la taille d'échantillon}
On étudie un modèle d’espérance $\mu=1$ et d’écart type $\sigma=2$.\begin{questions}
\item Calculer le seuil $2\sigma/\sqrt n$ pour $n=16$ puis $n=64$.
\item Le critère est-il satisfait pour $m=2{,}2$, $n=16$ ? Et pour $m=1{,}4$, $n=64$ ?
\item Par quel facteur faut-il multiplier $n$ pour diviser le seuil par 3 ?
\end{questions}
\end{exercice}
\end{blocLivre}
\clearpage
\rubrique{Problèmes}
\begin{blocLivre}{ex-11-12}
\begin{exercice}[3]{Simuler une loi à trois valeurs}
$X$ prend les valeurs $-1$, 2 et 5 avec probabilités $0{,}5$, $0{,}3$ et $0{,}2$.\begin{questions}
\item Calculer $\mu$ et $\sigma$.
\item Écrire une fonction de simulation avec un seul appel à \code{random()}.
\item Écrire une fonction qui renvoie la moyenne d’un échantillon de taille $n\geq1$.
\item Expliquer comment comparer les moyennes pour $n=10$, 100 et 1000 sans confondre résultat observé et preuve.
\end{questions}
\end{exercice}
\end{blocLivre}
\begin{blocLivre}{ex-11-13}
\begin{exercice}[3]{Répéter cinq cents échantillons}
Pour une variable de Bernoulli de paramètre $0{,}4$, on simule 500 échantillons de taille 25.\begin{questions}
\item Calculer l’espérance, l’écart type et le seuil $2\sigma/\sqrt{25}$.
\item Proposer un programme qui compte les moyennes satisfaisant le critère et renvoie leur proportion.
\item Le résultat doit-il obligatoirement être 95\,\% ?
\end{questions}
\end{exercice}
\end{blocLivre}
\begin{blocLivre}{ex-11-14}
\begin{exercice}[3]{Comparer deux groupes et leur réunion}
Le groupe A donne les valeurs $4,5,5,6$, le groupe B les valeurs $1,5,5,9$.\begin{questions}
\item Comparer moyennes et variances descriptives.
\item Calculer la moyenne et la variance de la série réunissant les huit valeurs.
\item Donner la médiane et les quartiles de la série réunie.
\item Expliquer ce que la moyenne seule masque ici.
\end{questions}
\end{exercice}
\end{blocLivre}
\begin{blocLivre}{ex-11-15}
\begin{exercice}[3]{Fréquences des lettres d'un texte}
On utilise les textes sans espaces ni accents \code{BANANE} et \code{ANANAS}.\begin{questions}
\item Calculer les fréquences des lettres présentes dans chaque texte.
\item Proposer une liste des fréquences pour les 26 lettres, avec \code{texte.count(c)} et \code{len(texte)}.
\item On choisit uniformément une position de \code{BANANE} et on note $V$ l’indicatrice d’une voyelle. Donner sa loi, son espérance et sa variance.
\end{questions}
\end{exercice}
\end{blocLivre}
